\documentclass[11pt]{article}
\usepackage{amsmath, amssymb}
\usepackage{geometry}
\usepackage{booktabs}
\usepackage{bm}
\usepackage{xcolor}

\geometry{margin=2.5cm}

\newcommand{\dt}{\Delta t}
\newcommand{\keff}{k_\mathrm{eff}}
\newcommand{\Mtt}{M_{\theta\theta}}
\newcommand{\Kth}{\mathbf{K}_\theta}
\newcommand{\Mss}{\mathbf{M}_{ss}}

% colour for the single changed entry
\newcommand{\highlight}[1]{\textcolor{red}{#1}}

\title{Implicit Torsional Spring: Saddle-Point Formulation\\
       and Modification from Scheme~A to Scheme~C}
\author{}
\date{}

\begin{document}
\maketitle

% ============================================================
\section{Full saddle-point system}
% ============================================================

The code solves a saddle-point linear system at each time step via GMRES.
For a single swimmer with $N_f$ filaments, each with $N_s$ segments, the
unknowns are
\[
  \mathbf{x} =
  \begin{pmatrix}
    \mathbf{f}_s   \\[2pt]
    \mathbf{f}_b   \\[2pt]
    \dot\psi_1     \\[2pt]
    \dot\psi_2
  \end{pmatrix},
\]
where $\mathbf{f}_s\in\mathbb{R}^{3N_fN_s}$ are the segment constraint forces,
$\mathbf{f}_b\in\mathbb{R}^{3N_b}$ are the blob forces,
$\dot\psi_1$ is the beat-phase rate, and $\dot\psi_2$ is the
shape-rotation rate.

\subsection*{Physical meaning of each block}

\paragraph{Mobility row ($\mathbf{f}_s$-block).}
Segment velocities are kinematically prescribed as the beat-phase
velocity plus the shape-rotation velocity:
\[
  \Mss\,\mathbf{f}_s
  + \mathbf{M}_{sb}\,\mathbf{f}_b
  - \mathbf{K}_\phi\,\dot\psi_1
  - \Kth\,\dot\psi_2
  = \mathbf{v}_s^\mathrm{ref}.
\]
Here $\Mss$ is the RPY/FCM hydrodynamic mobility between segments,
$\mathbf{K}_\phi$ is the reference velocity field for a unit phase
increment (\texttt{vel\_dir\_phase}), and $\Kth$ is the reference
velocity field for a unit rotation increment (\texttt{vel\_dir\_angle}).
The constraint force $\mathbf{f}_s$ is whatever internal (elastic +
inextensibility) force is required to maintain the prescribed velocity.

\paragraph{Phase row ($\dot\psi_1$-block).}
Generalised force balance projected onto the phase mode.
Only segment forces appear here; blob forces do not couple to the phase
or rotation DOFs:
\[
  -\mathbf{K}_\phi^\top\mathbf{f}_s = -q_\phi.
\]

\paragraph{Rotation row ($\dot\psi_2$-block).}
Generalised force balance projected onto the rotation mode, with
torsional spring restoring torque $\keff\psi_2^n$:
\[
  -\Kth^\top\mathbf{f}_s = -(q_\theta - \keff\psi_2^n).
\]

% ============================================================
\section{Scheme A (explicit spring)}
% ============================================================

In Scheme~A the spring is evaluated at the known level~$n$ and placed
entirely on the right-hand side; the matrix is unchanged:
%
\begin{equation}\label{eq:schemeA}
\underbrace{
\begin{pmatrix}
  \Mss & \mathbf{M}_{sb} & -\mathbf{K}_\phi & -\Kth \\[3pt]
  \mathbf{M}_{bs} & \mathbf{M}_{bb} & \mathbf{0} & \mathbf{0} \\[3pt]
  -\mathbf{K}_\phi^\top & \mathbf{0}^\top & 0 & 0 \\[3pt]
  -\Kth^\top & \mathbf{0}^\top & 0 & 0
\end{pmatrix}
}_{\mathbf{A}^{(\mathrm{A})}}
\begin{pmatrix}\mathbf{f}_s\\\mathbf{f}_b\\\dot\psi_1\\\dot\psi_2\end{pmatrix}
=
\underbrace{
\begin{pmatrix}
  \mathbf{v}_s^\mathrm{ref}\\[3pt]
  \mathbf{v}_b^\mathrm{ref}\\[3pt]
  -q_\phi\\[3pt]
  -(q_\theta-\keff\psi_2^n)
\end{pmatrix}
}_{\mathbf{b}}.
\end{equation}
%
The $(\dot\psi_2,\dot\psi_2)$ diagonal entry is \textbf{zero}.

\subsection*{Derivation of the amplification factor}

To isolate the update rule for $\psi_2$, work with the rotation mode
only (decouple from the phase mode and blobs by setting
$\mathbf{v}_b^\mathrm{ref}=\mathbf{0}$, $q_\phi=0$, $\mathbf{f}_b=0$,
$\dot\psi_1=0$) and consider the homogeneous problem $q_\theta=0$.
The system reduces to
\begin{equation}\label{eq:reduced_A}
\begin{pmatrix} \Mss & -\Kth \\ -\Kth^\top & 0 \end{pmatrix}
\begin{pmatrix} \mathbf{f}_s \\ \dot\psi_2 \end{pmatrix}
=
\begin{pmatrix} \mathbf{0} \\ \keff\psi_2^n \end{pmatrix}.
\end{equation}
Eliminate $\mathbf{f}_s$ by solving the first block row:
\begin{equation}\label{eq:fs_A}
  \Mss\,\mathbf{f}_s = \Kth\,\dot\psi_2
  \quad\Longrightarrow\quad
  \mathbf{f}_s = \Mss^{-1}\Kth\,\dot\psi_2.
\end{equation}
Substitute into the second block row:
\begin{equation}\label{eq:schur_A}
  -\Kth^\top \Mss^{-1}\Kth\,\dot\psi_2 = \keff\psi_2^n.
\end{equation}
The scalar $\Xi \equiv \Kth^\top\Mss^{-1}\Kth > 0$ is the Schur complement
(positive because $\Mss$ is SPD).  Dividing both sides by $-\Xi$:
\begin{equation}\label{eq:phidot_A}
  \dot\psi_2 = -\frac{\keff}{\Xi}\,\psi_2^n.
\end{equation}
Define the \emph{effective rotation-mode mobility}
\begin{equation}\label{eq:Mtt}
  \boxed{\Mtt \;\equiv\; \frac{1}{\Xi}
    = \bigl(\Kth^\top\Mss^{-1}\Kth\bigr)^{-1} > 0.}
\end{equation}
Then
\begin{equation}\label{eq:dotpsi_A}
  \dot\psi_2 = -\keff\,\Mtt\,\psi_2^n,
\end{equation}
and the explicit Euler update $\psi_2^{n+1}=\psi_2^n+\dt\,\dot\psi_2$ gives
\begin{equation}\label{eq:amp_A}
  \boxed{\psi_2^{n+1} = \bigl(1 - \keff\,\Mtt\,\dt\bigr)\,\psi_2^n.}
\end{equation}
The amplification factor is $\lambda_A = 1 - \keff\Mtt\dt$.
This lies in $(-1,1)$ only when $0 < \keff\Mtt\dt < 2$, i.e.\
\[
  \dt < \frac{2}{\keff\,\Mtt}.
\]
For $\mathtt{spring\_factor}=1$, $L=52$, $\Mtt\sim O(1)$, one obtains
$\keff\approx8.8\times10^5$, requiring $\dt\lesssim2\times10^{-6}$.
The constraint becomes prohibitive long before $\keff\Mtt\dt$ reaches~2.

% ============================================================
\section{Scheme C (backward-Euler spring)}
% ============================================================

\subsection*{Why can an entry be added to the saddle-point matrix?}

The saddle-point system is simply the equations of motion written in
matrix form.  There is nothing special about the structure that forbids
adding new entries --- what matters is where each physical term lands
after discretisation.

The bottom row of the system is a \emph{torque balance} on the rotation
mode:
\[
  \text{(hydrodynamic torque on }\dot\psi_2\text{)}
  = \text{driving torque} - \text{spring torque}.
\]
In Scheme~A the spring torque is evaluated at the \emph{old} time level
$\psi_2^n$, which is a known number.  It therefore contributes only to
the right-hand side, leaving the $(\dot\psi_2,\dot\psi_2)$ diagonal
entry of the matrix as zero.

In Scheme~C the spring is evaluated at the \emph{new} time level
$\psi_2^{n+1}$.  Because $\psi_2^{n+1} = \psi_2^n + \dt\,\dot\psi_2$,
the spring torque splits into two parts:
\[
  \keff\,\psi_2^{n+1}
  = \underbrace{\keff\,\psi_2^n}_{\text{known} \Rightarrow \text{RHS}}
  + \underbrace{\keff\dt\,\dot\psi_2}_{\text{unknown} \Rightarrow \text{LHS}}.
\]
The first part is known and stays on the right-hand side unchanged.
The second part depends on the unknown $\dot\psi_2$, so it must move
to the left-hand side, where it becomes a diagonal entry in the matrix.
That is the $-\keff\dt$ entry (the minus sign comes from the sign
convention of the saddle-point system).

In other words, the new matrix entry is not added arbitrarily: it arises
directly from choosing to evaluate the spring torque at the new time
level.  The right-hand side is identical to Scheme~A because the known
part of the spring, $\keff\psi_2^n$, is unchanged.

\subsection*{Resulting system}

Scheme~C treats the spring implicitly.  Replacing $\psi_2^{n+1}$ with
$\psi_2^n+\dt\,\dot\psi_2$ in the spring force:
\[
  \keff\psi_2^{n+1} = \keff\bigl(\psi_2^n + \dt\,\dot\psi_2\bigr).
\]
Moving the new $\keff\dt\,\dot\psi_2$ term to the left-hand side of the
rotation row changes exactly one matrix entry:
\[
  -\Kth^\top\mathbf{f}_s - \keff\dt\,\dot\psi_2
  = -(q_\theta - \keff\psi_2^n),
\]
so the $(\dot\psi_2,\dot\psi_2)$ diagonal changes from $0$ to $-\keff\dt$:
%
\begin{equation}\label{eq:schemeC}
\underbrace{
\begin{pmatrix}
  \Mss & \mathbf{M}_{sb} & -\mathbf{K}_\phi & -\Kth \\[3pt]
  \mathbf{M}_{bs} & \mathbf{M}_{bb} & \mathbf{0} & \mathbf{0} \\[3pt]
  -\mathbf{K}_\phi^\top & \mathbf{0}^\top & 0 & 0 \\[3pt]
  -\Kth^\top & \mathbf{0}^\top & 0 & \highlight{-\keff\dt}
\end{pmatrix}
}_{\mathbf{A}^{(\mathrm{C})}}
\begin{pmatrix}\mathbf{f}_s\\\mathbf{f}_b\\\dot\psi_1\\\dot\psi_2\end{pmatrix}
=
\mathbf{b}.
\end{equation}
%
The right-hand side $\mathbf{b}$ is \emph{identical} to Scheme~A.
The modified entry is highlighted.

\subsection*{Derivation of the amplification factor}

Under the same reduction as before, \eqref{eq:reduced_A} becomes
\begin{equation}\label{eq:reduced_C}
\begin{pmatrix} \Mss & -\Kth \\ -\Kth^\top & -\keff\dt \end{pmatrix}
\begin{pmatrix} \mathbf{f}_s \\ \dot\psi_2 \end{pmatrix}
=
\begin{pmatrix} \mathbf{0} \\ \keff\psi_2^n \end{pmatrix}.
\end{equation}
The first block row is unchanged, so \eqref{eq:fs_A} still holds:
$\mathbf{f}_s = \Mss^{-1}\Kth\,\dot\psi_2$.
Substituting into the second block row (which now has the extra
$-\keff\dt$ term):
\begin{equation}\label{eq:schur_C}
  -\Kth^\top\Mss^{-1}\Kth\,\dot\psi_2 - \keff\dt\,\dot\psi_2
  = \keff\psi_2^n.
\end{equation}
Collecting the $\dot\psi_2$ terms, using $\Xi = 1/\Mtt$:
\begin{equation}\label{eq:collect_C}
  -\bigl(\Xi + \keff\dt\bigr)\,\dot\psi_2 = \keff\psi_2^n
  \quad\Longrightarrow\quad
  \dot\psi_2 = -\frac{\keff\psi_2^n}{\Xi + \keff\dt}.
\end{equation}
The update $\psi_2^{n+1} = \psi_2^n + \dt\,\dot\psi_2$ gives
\begin{equation}\label{eq:amp_C_Xi}
  \psi_2^{n+1}
  = \psi_2^n - \frac{\keff\dt}{\Xi+\keff\dt}\,\psi_2^n
  = \frac{\Xi}{\Xi+\keff\dt}\,\psi_2^n.
\end{equation}
Expressing the result in terms of $\Mtt = 1/\Xi$,
divide numerator and denominator by $\Xi$:
\begin{equation}\label{eq:amp_C}
  \boxed{\psi_2^{n+1}
  = \frac{1}{1 + \keff\,\Mtt\,\dt}\,\psi_2^n.}
\end{equation}
The amplification factor is $\lambda_C = 1/(1+\keff\Mtt\dt)$.
\begin{itemize}
  \item \textbf{Unconditionally stable}: for all $\keff>0$, $\dt>0$,
        $\Mtt>0$ one has $0 < \lambda_C < 1$, so $\psi_2^n\to0$
        monotonically.
  \item \textbf{No time-step restriction}: $\lambda_C\to0$ as
        $\keff\dt\to\infty$ rather than diverging.
\end{itemize}

Comparing the two factors:
\begin{center}
\begin{tabular}{@{}lll@{}}
\toprule
 & Scheme A & Scheme C \\
\midrule
Schur complement for $\dot\psi_2$ & $-\Xi$ & $-(\Xi + \keff\dt)$ \\
$\dot\psi_2$ (homogeneous) & $-\keff\Mtt\psi_2^n$ &
  $\displaystyle-\frac{\keff\Mtt}{1+\keff\Mtt\dt}\psi_2^n$ \\[6pt]
Amplification $\lambda$ & $1 - \keff\Mtt\dt$
  & $\dfrac{1}{1+\keff\Mtt\dt}$ \\[6pt]
Stable when & $\keff\Mtt\dt < 2$ & always \\
\bottomrule
\end{tabular}
\end{center}

\paragraph{No-rotation limit.}
As $\keff\to\infty$ the augmented Schur complement
$\Xi+\keff\dt\to\infty$, so $\dot\psi_2\to0$ and $\psi_2^n\to0$.
The coupling term $\Kth\dot\psi_2\to\mathbf{0}$ decouples the phase
mode from the rotation mode, and the effective phase-mode mobility
converges exactly to $M_{\phi\phi}=(\mathbf{K}_\phi^\top\Mss^{-1}\mathbf{K}_\phi)^{-1}$
--- the result obtained with rotation disabled.  In Scheme~A the matrix
Schur complement is independent of $\keff$, so this exact limit is
never reached.

% ============================================================
\section{Code implementation}
% ============================================================

All changes are in \texttt{src/cilia/mobility/mobility\_solver.cpp}.
\texttt{filament.cpp} is unchanged.

\subsection*{Matrix--vector product (\texttt{system\_matrix\_mult},
             \texttt{system\_matrix\_mult\_new})}

The existing code computes $-\Kth^\top\mathbf{f}_s$ via a
$k$-loop over segments.  Immediately after that loop, add the
$(\dot\psi_2,\dot\psi_2)$ diagonal:

\begin{verbatim}
#if DYNAMIC_SHAPE_ROTATION
{
  const Real k_eff_dt = TORSIONAL_SPRING_MAGNITUDE_FACTOR
                        * 2.0*M_PI*FIL_LENGTH*FIL_LENGTH*FIL_LENGTH*DT;
  out(in_pos_theta) -= k_eff_dt * shape_rotation_angle_dot;
}
#endif
\end{verbatim}

\noindent where \texttt{in\_pos\_theta} is the index of $\dot\psi_2$ in
the solution vector (\texttt{in\_pos + NSWIM*NFIL} in the
\texttt{DYNAMIC\_PHASE\_EVOLUTION} branch; \texttt{in\_pos} in the
rotation-only branch).  The minus sign reflects the saddle-point sign
convention.

\subsection*{Preconditioner (\texttt{apply\_preconditioner})}

The simple diagonal preconditioner for $\dot\psi_2$ inverts the
diagonal approximation of $-\Xi$, namely $-\|\Kth\|^2/\mu_\mathrm{seg}$.
With Scheme~C the effective diagonal is $-(\Xi + \keff\dt)$, so the
divisor augments from $\|\Kth\|^2/\mu_\mathrm{seg}$ to:
\begin{verbatim}
out(...) /= Knorm_angle * seg_mob_inv + k_eff_dt;
\end{verbatim}

\subsection*{Coupled $2\times2$ preconditioner
             (\texttt{apply\_preconditioner\_new})}

For the $(\dot\psi_1,\dot\psi_2)$ coupled block the preconditioner
inverts a $2\times2$ matrix $C$ with $c_{22}\approx\|\Kth\|^2$.
The mobility scaling $\mu_\mathrm{seg}$ is applied after inversion, so
augmenting the $(\dot\psi_2,\dot\psi_2)$ entry by $-\keff\dt$ in the
system matrix corresponds to
\begin{verbatim}
c22 += k_eff_dt / seg_mob_inv;
\end{verbatim}
before the $2\times2$ matrix inverse.

\subsection*{Right-hand side}

No change.  The term $-(q_\theta-\keff\psi_2^n)$ on the last row of
$\mathbf{b}$ was already present in Scheme~A and is retained unchanged
in Scheme~C.

% ============================================================
\section{Summary}
% ============================================================

\begin{center}
\begin{tabular}{@{}lll@{}}
\toprule
 & Scheme A & Scheme C \\
\midrule
$(\dot\psi_2,\dot\psi_2)$ matrix entry  & $0$
  & $-\keff\dt$ \\
RHS last row  & $-(q_\theta-\keff\psi_2^n)$ & same \\
Schur complement for $\dot\psi_2$  & $-\Xi$
  & $-(\Xi+\keff\dt)$ \\
Amplification factor $\lambda$ & $1-\keff\Mtt\dt$
  & $1/(1+\keff\Mtt\dt)$ \\
Stability condition  & $\keff\Mtt\dt < 2$ & none \\
$\keff\to\infty$ limit  & approximate & exact no-rotation \\
\bottomrule
\end{tabular}
\end{center}

\end{document}
